\section{Environment and observation conditions}
\label{sec:setup}

Four players start with \$1,500 on the forty-space United States board. The bank has 32 houses and twelve hotels. Declining a purchase opens an ascending auction. Bilateral trade threads permit proposals and counteroffers before ordinary play resumes. A color group is complete when one player owns every deed, including mortgaged deeds. Construction also requires cash, unmortgaged holdings, bank stock, and even building. The engine supplies the legal actions and alone applies state changes.

Debt management permits sales, mortgages, or bankruptcy, but no new rescue trade. A player may declare bankruptcy despite a payable debt. Player-creditors inherit buildings intact and receive mortgaged deeds without a transfer-interest charge. Assets returned to the bank are not immediately auctioned. Developed groups cannot be traded. Build plans face entry-by-entry evenness checks. Hotel downgrades appear on the liquidation menu only when four bank houses exist. These implementation conventions affect financing and succession. They are not all stated in the delivered system prompt; Appendix~\ref{app:conditions} specifies the information boundary and preserved arithmetic convention.

The cohort has twenty completed games, eighty player-games, and 50,918 events. Five seeds each supply four cyclic seat rotations. Every configuration occupies every absolute seat five times, but its clockwise neighbors remain fixed. Games run for 75 to 386 global turns and end through bankruptcy without a turn cap. Initial seed equality does not ensure equal later random exposure because actions change random-number consumption. G$n$ d$k$ identifies game $n$ and zero-based decision $k$.

\begin{table}[t]
\centering\small
\begin{tabularx}{\linewidth}{@{}lYrr@{}}\toprule
Label & OpenRouter model identifier & Decisions & Attempts\\\midrule
Sol & \path{openai/gpt-6.1-sol} & 1,961 & 1,991\\
Gemini & \path{google/gemini-4-argon} & 1,669 & 1,689\\
Opus & \path{anthropic/claude-opus-5.5} & 1,696 & 3,301\\
Grok & \path{x-ai/grok-4.7} & 1,663 & 1,698\\\bottomrule
\end{tabularx}
\caption{The four historical configurations. A medium reasoning-effort request is recorded for every decision. Provider-reported token fields do not make computation comparable. All requests use OpenRouter; temperature and nucleus sampling parameters are not supplied.}
\end{table}

The standard prompt asks the player to win by any legal strategy and invites public bluffing and deception. It requires both a public message and an honest private note. Delivered states contain cash, holdings, buildings, bank stock, twenty public timeline entries, and ten own private notes. Purchase prompts include the landed property's economic parameters; management, trade, jail, and liquidation prompts identify deeds without a complete price, rent, mortgage, or redemption table. The system prompt contains no comprehensive rules reference. Missing current information does not establish missing model knowledge, but it makes prior knowledge and history retrieval part of the task. Full standard and recovery prompts are printed in Appendix~\ref{app:prompts}.

Opus receives first-refusal recovery in 1,604 of 1,696 decisions. All 1,605 refused responses carry the same restriction message about reverse engineering or duplicating model outputs. The first-attempt refusal counts are 64/84, 40/60, 30/53, and 80/109 in G1--G4, then 100\% in G5--G20. Non-refused first attempts stop during the first four games, with last responses between 16:15 and 16:38 UTC on 25 October. This is a within-campaign chronology, not evidence of a causal code-version change.

The recovery request removes private-note production and the request for honest private reasoning, adds a prohibition on returning a reasoning transcript, and retains permission to deceive publicly. It also consumes the remaining attempt. Current Anthropic documentation describes refusals of internal-reasoning extraction requests~\cite{anthropic2026opus}. That is consistent with a prompt/provider incompatibility; the records contain no exact classifier field establishing the trigger. Opus consequently lacks authored notes in 1,394 delivered memory windows and has less opportunity to correct invalid recovery actions. It still sees public history and current state. These conditions also change the market faced by its opponents.

Grok uses required tool choice. Among 1,068 selected responses with an extractable action name in assistant content, 139 name a different action from the received tool call. The engine applies the tool call. The disagreement appears in saved response bodies, with its upstream origin unresolved. Sixty-four of Grok's 136 thread openings and seven of its twenty-eight acceptances have such disagreements. Current trade terms are reconciled against legal acceptance menus, although 2,054 initial prompts contain defective past-trade summaries. We preserve actual exchanges and inspect delivered context when discussing what an opponent encountered. Appendix~\ref{app:conditions} records the remaining conditions and affected counts once.

\section{Measures and analytical scope}
\label{sec:methods}

An offer is one proposal or counter awaiting a response. The current author may differ from the thread initiator. We calculate ownership completion atomically and use the 918 offers whose menus permit acceptance, keeping the two ineligible offers in the overall response census. Opening offers and later offers are separated because continuation depends on earlier responses. Joint recipient-by-stage counts make the conditioning explicit. We do not treat offers in a thread or turns in a game as independent trials.

For pricing, we restrict the ledger to cash exchanged for deeds, excluding jail cards and countervailing assets or cash. An ask comes from the deed seller; a bid comes from the cash buyer. The ratio divides cash by the sum of deeds' board list prices. That denominator is a reproducible scale, not a fair or strategic value. We stratify by author, mortgage status, and thread position, then inspect single-deed offers within a game and seller. The latter comparisons hold fewer variables uncontrolled but do not fix timing, buyer identity, cash, outside options, or prior dialogue. Sparse overlap limits the claim rather than licensing a fitted completion coefficient.

A completion episode begins when an acquisition creates a complete group and ends at ownership loss or game end. For all sixty-three trade episodes, the first management window is the owner's first actual management turn after settlement, including its successive decisions. We retain prior classifications of construction access and execution. Fresh arithmetic adds acquisition and build turns, ownership duration, paid rent, and opponent movement events during the episode. Each opponent move is counted once per exposed group, including card-driven moves. A group landing is a recorded move ending on any deed in the group, even if mortgaged or nonpaying. This is realized exposure, not a probability model or an exogenous opportunity treatment.

Debt episodes join consecutive liquidation decisions for one player, turn, creditor, reason, and owed amount, ending at a non-liquidation decision or elimination. Each of the 128 successful episodes has a matching payment event in the final action's emitted range. The sixty elimination endpoints join the preserved elimination ledger. Existing legal payment sequences establish three payable eliminations. Fifty-six others have debt above the generous bound of cash plus all building salvage plus every unmortgaged deed's mortgage principal, including collateral unlocked by removing buildings. One has sufficient nominal assets but a menu-level hotel lock. No engine search or witness execution is added here.

First-estate positions use the simultaneous delivered state before the first bankruptcy action in each game. We compare cash, deeds, building house-equivalents, and paid rent previously received. A hotel counts as five construction levels for this accounting display, not five available bank houses. Rank one includes ties. These are transparent descriptions of prior strength, without a constructed utility score. They do not estimate the independent predictive value or causal effect of inheritance.

Semantic evidence remains inherited. Two Astra readers labeled the seventy-four outright rejections of acceptance-legal recipient-completing offers. Both assigned a rival-consequence code in sixty-one records; each assigned sixty-three positives, producing seventy agreements and four disagreements. The code covers explicit rival development, ownership, finance, survival, or estate consequences in applied public and private fields. The second reader had partial prior exposure, and neither is independent human validation. The eight incidents use source-exact fields, preserve channel and attempt identity, and distinguish recipient delivery from agreement with a proposition. Explanations can be unfaithful~\cite{turpin2023unfaithful}; their words establish what was recorded, not what internally caused the response.

All quantitative claims concern the finite historical cohort. No new games, evaluations, synthetic tests, baselines, engine runs, or replays are performed. New scripts read existing records and hash every numerical input. Historical engine witnesses and replay results retain their prior status. With five dependent seed blocks and no repeated seed-seat cell, bootstrap intervals would offer little reliable information about a new model, roster, or provider regime. We report complete cells, small overlap sets, and trajectory-level variation instead of independent-offer significance tests.
